\section{The arithmetic of normalized theta coefficients}
\label{sec:theta-arithmetic}

The arithmetic premise retained in PA15 can now be supplied. We prove
the actual one-dimensional statements \texttt{galSg}, \texttt{gal1},
and \texttt{use1}, including their coefficient-field and torsion-point
conditions. This section uses the literal host and curve datum of PA,
not a column law, a vacuum identity, or a later-stage geometric point.
It does not yet prove the multidimensional vacuum or covariance
identities. The ordinary uniformization theorem for a nonsingular
complex elliptic curve is a named classical input here. We explain
its precise uses, and the intervening algebraic descent arguments,
instead of incorporating that whole classical theory as a new axiom
with the desired theta conclusion.

Write $\Lambda=\lambda\iota(I)$ for the curve lattice, with
$\Delta_\Lambda=1$, and $A=\operatorname{area}(\Lambda)>0$. Let
$s_\Lambda=\mathrm{sgStar}(\Lambda)$ be the normalized odd theta
function described in PA15. Its derivative at zero is one, its
zeros are exactly $\Lambda$, and
\[
 s_\Lambda(z+\omega)=\eta(\omega)
  \exp\!\left(\frac{\pi}{A}(z+\omega/2)\bar\omega\right)s_\Lambda(z),
 \quad
 \eta(\omega)=\begin{cases}1&\omega\in2\Lambda,\\-1&\omega\notin2\Lambda.
 \end{cases}                                                    \tag{TA.1}
\]
The translation operator is the actual operator
\[
 U_xs(z)=s(z+x)\exp\!\left(-\frac{\pi}{A}z\bar x
                           -\frac{\pi}{2A}x\bar x\right).       \tag{TA.2}
\]
All statements about coefficients below use $[z^n]f=f^{(n)}(0)/n!$.
There is no claim that a general field embedding commutes with limits,
complex differentiation, or the exponential function.

\paragraph{TA14: analytic normalization without an arithmetic premise.}
We isolate this preliminary claim so that the dependency graph does
not use the arithmetic assumption of PA15 to prove that assumption.
For a genuine lattice the canonical product
\[
 \sigma(z)=z\prod_{0\ne\omega\in\Lambda}
       (1-z/\omega)\exp(z/\omega+z^2/(2\omega^2))
\]
converges normally on compact sets: the logarithm of each factor
is $O(|z/\omega|^3)$ away from finitely many factors, and the
lattice sum of $|\omega|^{-3}$ converges. It has simple zeros
exactly at $\Lambda$, is odd by pairing $\omega,-\omega$, and
has derivative one at zero. Differentiating its logarithm
gives $(\log\sigma)''=-\wp$. Therefore the increments of its
logarithmic derivative along periods are constants
$\eta_\omega$, additive in $\omega$. Integrating around a
positively oriented fundamental parallelogram, with the one
residue of the logarithmic derivative, gives Legendre's
identity $\eta_1\omega_2-\eta_2\omega_1=2\pi i$.
For $\operatorname{Im}(\bar\omega_1\omega_2)=A>0$, this says
$\eta_\omega=c\omega+\pi\bar\omega/A$ for one $c\in\mathbb C$.
The sign in the integrated period law is $-1$ on each basis
period, by oddness at its half-period. Composing the two laws
and using Legendre gives the sign
$(-1)^{r+s+rs}$ at $r\omega_1+s\omega_2$, which is one
exactly when both integers are even. Hence
$f(z)=e^{-cz^2/2}\sigma(z)$ has precisely (TA.1), with its
stated zeros, parity and derivative.

If an entire $G$ has those same translation laws, $G/f$ is
elliptic and has at most a single simple pole on the torus.
The residue theorem on a fundamental parallelogram forces
that residue to vanish, so it has no pole. The maximum
principle makes it constant; its value is $G'(0)$.
This proves uniqueness and identifies the source's chosen
$\mathrm{sgStar}$ with $f$, without any field-membership
assertion. The same uniqueness gives the scaling identity
$s_{a\Lambda}(az)=a s_\Lambda(z)$ for $a\ne0$, since
$\operatorname{area}(a\Lambda)=|a|^2A$.

Finally, for $v\notin\Lambda$ the function
$-f(u+v)f(u-v)/(f(u)^2f(v)^2)$ is elliptic in $u$,
has just the double pole $u^{-2}$ at zero, and vanishes
at $u=v$. Subtracting $\wp(u)$ and cancelling poles shows
that it is $\wp(u)-\wp(v)$. This is the exact sigma addition
identity used in TA3 and TA7. All analytic facts used below
are therefore available independently of odd-theta
coefficient algebraicity.

\paragraph{TA1: the actual conjugate subgroup is a period lattice.}
The definition of $\mathrm{lamConj}(M,\rho,\Lambda)$ selects an
additive subgroup with invariants $\rho(g_2),\rho(g_3)$. One must
justify that this is the correct lattice, rather than assume that
an arbitrary chosen subgroup is discrete. Here is the missing
classification argument. If either Eisenstein sum of orders four
or six is nonzero, it is summable: in the source the totalized sum
of a nonsummable family is zero. Summability implies that only
finitely many nonzero subgroup elements can have absolute value
at most one. Otherwise the corresponding inverse powers would
not tend to zero on the cofinite filter. Removing those finitely
many points leaves a neighborhood of zero containing only zero.
Thus the subgroup is discrete.

A discrete subgroup of the real plane has rank at most two.
For completeness, choose a shortest nonzero vector; if the group
is not contained in its real line, reduce along that vector and
choose a least positive perpendicular height. Subtracting integral
multiples of the two vectors places any point in a bounded
fundamental parallelogram. Discreteness and the minimal choices
then give an integral basis. The one-dimensional case is cyclic.
For a cyclic subgroup $a\mathbb Z$ the absolutely convergent sums give
\[
 g_2=\frac{4\pi^4}{3a^4},\qquad
 g_3=\frac{8\pi^6}{27a^6},\qquad g_2^3-27g_3^2=0;
\]
the zero subgroup also has discriminant zero. Hence a subgroup
with nonzero discriminant must be a rank-two lattice.

For two rank-two lattices with the same $g_2,g_3$, their Weierstrass
functions have the same Laurent series at zero. One can obtain
uniqueness directly from the coefficient comparison in TA2 below,
applied to $z^2\wp(z)$. The identity theorem extends equality as
meromorphic functions on $\mathbb C$, and their pole sets are the
two lattices. Thus the lattices are equal.

Existence for every pair $a,b$ with $a^3-27b^2\ne0$ follows from
classical elliptic uniformization: the smooth projective cubic
$y^2=x^3-ax/4-b/4$, with differential $dx/(2y)$, is
$\mathbb C/\Lambda'$ with $x=\wp_{\Lambda'}$ and
$y=\wp'_{\Lambda'}/2$. The differential normalization fixes the
scale, so its invariants are exactly $a,b$. Applying this with
$a=\rho(g_2)$ and $b=\rho(g_3)$ is legitimate because $\rho$ is
injective and the discriminant is one. Consequently the selected
\texttt{lamConj} is this unique genuine lattice $\Lambda'$.
In particular, if $\rho$ fixes both invariants, $\Lambda'=\Lambda$.
This also supplies the arbitrary-subgroup uniqueness clause of
the invariant dictionary, which PA1 deliberately did not assume.

\paragraph{TA2: a formal uniqueness theorem with all initial data.}
For a formal series $F$ over a characteristic-zero field set
\[
 W_s(F)=F'^2-FF''-sF^2.
\]
The cleared Weierstrass equation is
\[
 (W_s(F)'F-2W_s(F)F')^2
     =4W_s(F)^3-aW_s(F)F^4-bF^6.                               \tag{TA.3}
\]
For fixed $s,a,b$ there is at most one solution with
$F(0)=0$, $[z]F=1$, and $[z^2]F=0$. To prove this, write $F=zV$,
$V(0)=1$, and put $P=W_s(F)V^{-2}$. Cancelling $V^6$ in (TA.3)
gives
\[
 (zP'-2P)^2=4P^3-az^4P-bz^6,\qquad P(0)=1.                  \tag{TA.4}
\]
If two solutions $P,Q$ first differ in degree $n$, write
$P-Q=z^nE$, $E(0)\ne0$. Subtract (TA.4), cancel $z^n$,
and take constant terms. The two sides are
$-4(n-2)E(0)$ and $12E(0)$. Thus $(n+1)E(0)=0$, a
contradiction. This proves uniqueness of $P$.

For two original solutions $F=zV$ and $G=zT$, equality of $P$
gives equality of the derivatives of $V'/V$ and $T'/T$.
Their constant terms vanish because $[z^2]F=[z^2]G=0$.
In characteristic zero a formal series with derivative zero is
constant, so $V'/V=T'/T$; consequently $(V/T)'=0$ and
$V/T=1$. This proves $F=G$. The degree-two coefficient of
(TA.3), using the three initial conditions, also gives
\[
                         s=-2[z^3]F.                         \tag{TA.5}
\]
These are purely formal operations. In particular they commute
with every field embedding, with no convergence hypothesis on
the coefficientwise image series.

\paragraph{TA3: the differential equation and a CM polynomial.}
The purely analytic canonical product construction in TA14 gives
$s_\Lambda(z)=e^{-cz^2/2}\sigma_\Lambda(z)$ for the ordinary
Weierstrass sigma function. Since
$(\log\sigma_\Lambda)''=-\wp_\Lambda$, one obtains
\[
 W_c(s_\Lambda)=\wp_\Lambda s_\Lambda^2.                     \tag{TA.6}
\]
Equivalently (TA.6) follows by differentiating the sigma addition
formula twice at zero, so it uses no arithmetic property of $c$.
Together with $\wp'^2=4\wp^3-g_2\wp-g_3$, it gives (TA.3)
for the Taylor series of $s_\Lambda$. Formula (TA.5) identifies
the constant $c$ from that series.

Let $\alpha=\iota(\sqrt{-d})$, so
$\alpha\bar\alpha=d=:N$, $\alpha\Lambda\subset\Lambda$,
$\alpha\ne0$, and $\alpha^2=-d\ne d$. By (TA.1),
\[
 R(z)=\frac{s_\Lambda(\alpha z)^2}{s_\Lambda(z)^{2N}}         \tag{TA.7}
\]
is even and $\Lambda$-periodic: its exponential multipliers
cancel because $\alpha\bar\alpha=N$, and its signs have been
squared. Its only poles modulo $\Lambda$ occur at zero, with
order at most $2N-2$. An even elliptic function with poles only
there is a polynomial in $\wp$: successively subtract a multiple
of $\wp^j$ to cancel the highest even pole, then apply the
maximum principle to the holomorphic remainder on the compact
torus. Thus $R=Q(\wp)$ for a polynomial $Q$.

Suppose now $M\subset\mathbb C$ contains $\alpha,g_2,g_3$ and
all Taylor coefficients of $s_\Lambda$. By (TA.5), $c\in M$;
by (TA.6), the Laurent series of $\wp$ lies in $M((z))$.
The Laurent series of (TA.7) also lies in $M((z))$. Since
$\wp^j=z^{-2j}+\cdots$, the same descending pole cancellation
shows that every coefficient of $Q$ belongs to $M$.
This supplies the polynomial descent required before applying
$\rho$; it is not an assumption that a complex polynomial
appearing in an elliptic-function argument is already defined
over the coefficient field. The source uses an $M$-linear
retraction for this descent; distinct pole orders give the
shorter argument here.

\paragraph{TA4: CM eliminates the possible quadratic exponential.}
Let $\rho:M\hookrightarrow\mathbb C$ fix $\iota(K)$, with
$g_2,g_3$ and all coefficients of $s_\Lambda$ in $M$.
Let $\Lambda'$ be TA1's conjugate lattice and $s'=s_{\Lambda'}$.
Write $c'$ for its constant in (TA.6). Multiplication of a
function by $e^{uz^2}$ changes its constant from $c'$ to
$c'-2u$; direct differentiation verifies
\[
 W_{c'-2u}(e^{uz^2}s')=e^{2uz^2}W_{c'}(s').                  \tag{TA.8}
\]
Take $u=(c'-\rho(c))/2$. The coefficientwise image under $\rho$
of the original Taylor series and the actual Taylor series of
$e^{uz^2}s'$ have identical initial data and satisfy (TA.3) with
the same $\rho(c),\rho(g_2),\rho(g_3)$. TA2 makes the series
equal. In particular the coefficientwise image has now been
identified with a convergent series by a formal uniqueness
argument, rather than by an illicit continuity assertion.

Clear the denominator in (TA.7), substitute (TA.6), and apply
$\rho$ coefficientwise. TA3 shows that every coefficient being
transported is in $M$. Taylor injectivity for entire functions
then gives the analytic identity
\[
 s'(\alpha z)^2e^{\kappa z^2}
       =s'(z)^{2N}\rho(Q)(\wp_{\Lambda'}(z)),\qquad
             \kappa=2u(\alpha^2-N).                          \tag{TA.9}
\]
The identity is first read off the lattice and then extended
after clearing all poles. At a period $\omega\in\Lambda'$,
the right side vanishes: its possible pole has order at most
$2N-2$, while $s'^{2N}$ has order $2N$. Hence
$\alpha\omega\in\Lambda'$. The squared period law for $s'$
can therefore be used on both sides of (TA.9). Translating by
$\omega$ and cancelling on the nonempty open set where the
factors do not vanish gives
\[
                  e^{\kappa(z+\omega)^2}=e^{\kappa z^2}.
\]
Both sides are entire, so this holds for every $z$. For a
nonzero period, differentiation gives $2\kappa\omega=0$;
thus $\kappa=0$, and $\alpha^2\ne N$ forces $u=0$.
We have proved the exact statement \texttt{galSg}:
\[
 \rho([z^n]s_\Lambda)=[z^n]s_{\Lambda'}\quad(n\ge0),          \tag{TA.10}
\]
under its original assumption that all coefficients on the
left belong to $M$. The next arguments will supply algebraicity;
it was not silently assumed in this step.

\paragraph{TA5: torsion-point transport and the exact CM action.}
Under Weierstrass uniformization the class of $t$ is the origin
if $t\in\Lambda$, and otherwise is
$(\wp_\Lambda(t),\wp'_\Lambda(t)/2)$ on
$y^2=x^3-g_2x/4-g_3/4$. The source relation \texttt{StarPt}
says precisely that the second point is the coordinatewise
$\rho$-image of the first, with both coordinates in $M$ for
a point other than the origin, or that both points are the origin.
The ordinate can be zero, in particular at nonzero $2$-torsion.
This interpretation
proves uniqueness modulo $\Lambda'$, invariance under replacing
$t$ by a period translate, compatibility with addition, and
preservation of the exact annihilating integers. The group law
uses rational chord-and-tangent formulas, including the origin
and vertical-tangent cases, and therefore commutes with field
embeddings. Existence follows by uniformizing the image point.

We also need the action of every $a\in\mathcal O_K$, not just
integer multiplication. The map $z\mapsto\iota(a)z$ on
$\mathbb C/\Lambda$ is an algebraic endomorphism of its cubic.
Indeed its coordinate functions are elliptic meromorphic
functions, hence rational expressions in $\wp,\wp'$; the
identities giving a morphism and its group law are algebraic.
For a complex field automorphism $\sigma$, conjugate these
expressions and their identities. The resulting morphism on
the conjugate cubic is holomorphic; its lift through the
covering $\mathbb C\to\mathbb C/\Lambda'$ is $z\mapsto\beta z$.
The pullback of $dx/(2y)$ shows
$\beta=\sigma(\iota(a))$: this is an algebraic identity of
differentials, so conjugation of its coefficients is legitimate.
If $\sigma$ fixes $\iota(a)$, the transported action is exactly
the same scalar, with no unspecified isogeny or change of scale.

The output coordinates of this CM action must also belong to
$M$. Apply the preceding calculation to every automorphism
fixing all of $M$. It fixes the invariants, the input point,
and $\iota(a)$; hence it fixes the output point. The fixed
field of $\operatorname{Aut}(\mathbb C/M)$ is $M$: an algebraic
element outside $M$ has a distinct conjugate in characteristic
zero, and an element transcendental over $M$ can be shifted
in a transcendence basis; either embedding extends to the
algebraically closed field $\mathbb C$. Thus each nonzero
output coordinate is in $M$. This supplies the field-membership
part of the literal \texttt{StarPt} conclusion as well as its
coordinate identity.

For an arbitrary embedding $\rho:M\to\mathbb C$ it suffices
to extend its restriction on the countable field generated by
the invariants, both the input coordinates and the output
CM coordinates just proved to belong to $M$, and the
specified CM scalars. Such an embedding extends to an
automorphism of $\mathbb C$: extend it first to algebraic
closures, choose transcendence bases of the complements
(both have cardinality $|\mathbb R|$), identify the bases,
and extend through algebraic closures again. We do not assert
extension of an arbitrary embedding on an uncountable $M$.
All quantities whose images matter were included in the
countable subfield, so the preceding calculation proves the
literal CM clause for $\rho$.

If every nonzero $n$-torsion point has its two coordinates in
$M$, use one such extension simultaneously on their finite
set. There are exactly $n^2$ points, identified with
$\Lambda/n\Lambda$; their images are distinct $n$-torsion
points and exhaust the target's $n^2$ points. This proves the
surjectivity clause with its exact extra hypothesis. We have
now accounted for all seven clauses of \texttt{galStar}.
In particular $\Lambda'$ is stable under $\iota(\mathcal O_K)$:
apply the CM clause to the origin, represented by $0$ and any
$\omega'\in\Lambda'$.

\paragraph{TA6: conjugate curve ideals with the required coprimality.}
For later use we supply \texttt{galCurve}, including its
arbitrary prescribed positive integer $n$. Let
$\Lambda'=\mathbb Z\omega_1+\mathbb Z\omega_2$ be the lattice
of TA1, stable under $\mathcal O_K$ by TA5. For a nonreal
integral $a\in K$, write
$\iota(a)\omega_1=r\omega_1+s\omega_2$ with $r,s\in\mathbb Z$.
Here $s\ne0$, since otherwise $\iota(a)=r$. Put
$\delta=\omega_1/s$ and
$I=\{b\in\mathcal O_K:\delta\iota(b)\in\Lambda'\}$.
Stability makes $I$ an ideal, and $s\in I$ makes it nonzero.
Both periods are $\delta$ times integral elements, namely
$s$ and $a-r$, so $\Lambda'=\delta\iota(I)$.

Every nonzero ideal class of a Dedekind domain has a
representative coprime to a given nonzero ideal $J$.
Here is the elementary form used. Choose $a\in I$ whose
valuation at each prime dividing $J$ is exactly that of $I$,
by the finite Chinese remainder theorem in $I/\mathfrak p I$.
Then $(a)=I A$ with $A+J=\mathcal O_K$. Choose $a'\in A$
with the same exact valuations at the primes of $J$ and write
$(a')=A B$, again with $B+J=\mathcal O_K$.
It follows that $(a)B=(a')I$, and
\[
 \Lambda'=\delta\,\iota(a/a')\,\iota(B).                     \tag{TA.11}
\]
Apply this with $J=(2dn)$. Then $B$ has norm prime to $2dn$:
if a rational prime dividing $2dn$ divided its norm, some
prime above it would divide $B$ and $J$, a contradiction.
Its norm is therefore prime both to $2d$, as required of a
curve ideal, and to the prescribed $n$.

The discriminant of $\Lambda'$ is one, since its two invariants
are conjugates of those of $\Lambda$. The scaling identity in
PA1 therefore gives the exact twelfth-root equation for the
new factor $\delta'=\delta\iota(a/a')$. Thus $(B,\delta')$
is a literal curve datum, not merely a homothetic torus.
This proves \texttt{galCurve} without using theta coefficients.

\paragraph{TA7: division quotients are determined over the point field.}
For $u\notin\Lambda$ and $j\ge0$ define the analytic germ
\[
 \Psi_j(z)=\frac{s_\Lambda(j(z+u))}
                 {s_\Lambda(z+u)^{j^2}}.                    \tag{TA.12}
\]
The denominator is nonzero near zero. The sigma addition
formula and its derivative give
\[
 \Psi_0=0,\quad\Psi_1=1,\quad
 \Psi_2=-\wp'_\Lambda(z+u),\quad
 (j+1)^2\Psi_j\Psi_{j+2}
       =\Psi_{j+1}\Psi_{j+1}''-(\Psi_{j+1}')^2.              \tag{TA.13}
\]
For the recurrence, differentiate
$s(v+w)s(v-w)=\operatorname{wr}(s)(v)s(w)^2
-\operatorname{wr}(s)(w)s(v)^2$, where
$\operatorname{wr}(f)=ff''-f'^2$, and substitute
$v=(j+1)(z+u)$ and $w=z+u$. The product rule
$\operatorname{wr}(fg)=\operatorname{wr}(f)g^2+f^2\operatorname{wr}(g)$
and the corresponding power rule cancel the common powers
of $s(z+u)$ and give (TA.13).

The system $p'=q$, $q'=6p^2-g_2/2$, with
$p(0)=\wp(u)$ and $q(0)=\wp'(u)$, determines all derivatives
of $p,q$ recursively. Consequently their Taylor coefficients
belong to the field generated by these three data and commute
with embeddings of that field. Induct using (TA.13), starting
with $\Psi_1,\Psi_2$. At each step $\Psi_j$ for $j>0$ is a
nonzero germ, although its value at zero can vanish: otherwise
$s(j(z+u))$ would vanish on an open set and hence identically.
Division takes place in the Laurent-series field. If the
quotient is known analytically to be holomorphic, its Laurent
series has no negative terms and its Taylor coefficients lie
in the same field. This proves coefficient descent and
conjugacy for every $\Psi_j$, including at torsion points
where dividing numerical values at zero would be invalid.

We will also use a formal root observation. If $F(0)\ne0$,
all coefficients of $F^r$ and $F(0)$ lie in a characteristic-zero
field, and $r>0$, then all coefficients of $F$ lie there.
The coefficient of degree $k$ in $F^r$ is
$rF(0)^{r-1}[z^k]F$ plus a polynomial in earlier coefficients.
The same recursion proves conjugacy of two such roots when
their constant terms are paired. The constant term must be
specified; a power identity alone leaves a root of unity.

\paragraph{TA8: the constant that resolves the torsion root ambiguity.}
Fix $m>0$, $u\notin\Lambda$, and $m^2u\in\Lambda$.
Put $k=m^2-1$ and
\[
 K(z)=(U_us_\Lambda(z))^{4m^2},\qquad K(0)\ne0.
\]
Direct cancellation of the exponential factors in (TA.2)
gives
\begin{align}
 U_{2mu}s_\Lambda(2mz)&=K(z)\Psi_{2m}(z),                    \tag{TA.14}\\
 s_\Lambda(2m^2z)&=K(z)^{m^2}\Psi_{2m^2}(z).                \tag{TA.15}
\end{align}
In the second identity $2m^2u\in2\Lambda$, so (TA.1) makes
its translation operator the identity.
There is also the essential constant identity
\[
                         K(0)=-\frac{\Psi_k(0)}{\Psi_{k+2}(0)}.
                                                                  \tag{TA.16}
\]
Indeed $(k+1)u=m^2u$ is a period. Apply (TA.1) to the
points $u$ and $-u$ with that common period, and use oddness.
The quotient of the resulting values is
$-\exp(-2\pi m^2|u|^2/A)$. The powers of $s(u)$ in the
quotient of (TA.12) contribute $s(u)^{4m^2}$, giving
(TA.16). Its denominator is nonzero because
$(k+2)u\equiv u\pmod\Lambda$. Also $m=1$ is impossible
in this case, so $k>0$ and there is no exceptional zero
numerator. This is the specific root normalization required
by the source, not an arbitrary choice of an algebraic root.

\paragraph{TA9: the literal Galois identity for every rational translate.}
Take precisely the hypotheses of \texttt{gal1}: $\rho$ fixes
$\iota(K)$, the two lattice invariants belong to $M$, $m>0$,
$m(x/2)\in\Lambda$, and
\[
 \mathrm{StarPt}(M,\rho,\Lambda;x/(2m),x'/(2m));              \tag{TA.17}
\]
all coefficients of $G=U_xs_\Lambda$ are assumed in $M$.
Write $u=x/(2m)$, $u'=x'/(2m)$. TA5 gives
$m^2u'\in\Lambda'$. If $u\in\Lambda$, then $u'\in\Lambda'$
and both $x/2,x'/2$ are periods. The translations are the
identity by (TA.1), so TA4 applies directly.

Otherwise TA7 pairs all division quotients for $u,u'$ under
$\rho$. First apply its field-membership version to the
original lattice alone. Equation (TA.14), the assumed
coefficients of $G$, and Laurent division by the nonzero
$\Psi_{2m}$ show that the coefficients of $K$ lie in $M$.
Equation (TA.15) then shows that all coefficients of
$s_\Lambda$ lie in $M$, since rescaling the variable by the
nonzero integer $2m^2$ is invertible coefficientwise.
We can therefore use TA4; we have supplied, rather than
assumed, its additional coefficient-field premise.

Transport (TA.15), using TA4 and TA7, and cancel the nonzero
division-quotient germ. The series of $K^{m^2}$ and
$(K')^{m^2}$ are paired. Formula (TA.16) pairs their constant
terms $K(0),K'(0)$; TA7's root recursion pairs $K,K'$.
Finally (TA.14), followed by inverse integer rescaling,
gives the exact conclusion
\[
 \rho([z^n]U_xs_\Lambda)=[z^n]U_{x'}s_{\Lambda'}\quad(n\ge0).
                                                                  \tag{TA.18}
\]
This proves \texttt{galois\_translate\_sgStar}, with the same
integer $m$ and the literal $2m$-division condition. In
particular no root-of-unity factor has been lost.

\paragraph{TA10: algebraic coefficients for all actual torsion translates.}
PA14 supplies algebraic $g_2,g_3$ and algebraic coordinates
for every torsion point of the normalized CM lattice. Let
$x$ be any rational-lattice point and choose $r>0$ with
$rx\in\Lambda$. For every automorphism
$\sigma\in\operatorname{Aut}(\mathbb C/\overline{\mathbb Q})$
apply TA9 with $M=\mathbb C$, $\rho=\sigma$, $m=2r$, and
$x'=x$. Its hypotheses are all now justified: the coefficient
membership is automatic in $\mathbb C$, $\sigma$ fixes $K$
and the invariants, and the coordinates of $x/(4r)$ are
algebraic torsion coordinates. TA1 identifies the conjugate
lattice with the original lattice. Hence $\sigma$ fixes
every coefficient of $U_xs_\Lambda$.

An element of $\mathbb C$ fixed by all these automorphisms
belongs to $\overline{\mathbb Q}$. To see this directly, a
transcendental element can be included in a transcendence
basis over $\overline{\mathbb Q}$; replacing it by itself
plus one defines an automorphism of the rational-function
field, which extends through its algebraic closure to
$\mathbb C$. It does not fix that element. Thus
\[
 [z^n]U_xs_\Lambda\in\overline{\mathbb Q}
       \quad(x\in\mathbb Q\Lambda,\ n\ge0).                 \tag{TA.19}
\]
The functions are entire by (TA.2). This supplies the full
odd-theta premise retained in PA15. Applying PA15 now gives
algebraic Taylor coefficients for every rational translate
of the normalized even theta function, and the exact
nonzero algebraic lattice-shift multipliers. Its normalization
denominator is nonzero because the chosen half-point is
outside the lattice.

\paragraph{TA11: the distinguished half-point is unique modulo the lattice.}
Put $w=(1+i\sqrt d)/2$. The hypotheses give
$w(1-w)=(d+1)/4=2k$ for an integer $k$.
On $V=(\tfrac12\Lambda)/\Lambda\simeq\mathbb F_2^2$, the
endomorphisms $w$ and $1-w$ are complementary idempotents:
their product is zero and their sum is the identity.
If a class $t$ has both images nonzero, neither idempotent
has rank zero; consequently each has rank one. There is
exactly one element with nonzero components in both
one-dimensional summands. It follows that any two points
satisfying the literal three unkilled clauses differ by
a period. This proves the source's
\texttt{two\_torsion\_unique} without enumerating parity cases.

The normalized even function
$\vartheta_{\Lambda,t}=U_ts_\Lambda/U_ts_\Lambda(0)$ is unchanged
when $t$ is replaced by a period translate: the translation
composition law in PA15 gives the same nonzero scalar in
numerator and denominator. Thus it depends only on this
unique distinguished class.

\paragraph{TA12: even-theta conjugacy, including embeddings not fixing $K$.}
First let $\rho:\mathbb C\hookrightarrow\mathbb C$ fix
$\iota(K)$. TA5 transports the $16$-torsion point $t/8$
to some $u$ modulo $\Lambda'$. Apply TA9 with $x=t$,
$x'=8u$, and $m=4$: its condition $4(t/2)=2t\in\Lambda$
and its $8$-division point are exactly the ones just chosen.
Addition, CM multiplication, and exact-order preservation
in TA5 show that $8u$ satisfies the three unkilled clauses.
By TA6, $\Lambda'$ is an actual $\mathcal O_K$-stable curve
lattice. TA11 makes $8u$ congruent to any specified
unkilled point $t'$ on it. Divide (TA.18) by its degree-zero
identity. We obtain
\[
 \rho([z^n]\vartheta_{\Lambda,t})
                   =[z^n]\vartheta_{\Lambda',t'}.            \tag{TA.20}
\]

For ordinary complex conjugation, the canonical product,
the area identity $A_{\bar\Lambda}=A_\Lambda$, and (TA.2)
give
\[
 s_{\bar\Lambda}(z)=\overline{s_\Lambda(\bar z)},\qquad
 \vartheta_{\bar\Lambda,\bar t}(z)
                 =\overline{\vartheta_{\Lambda,t}(\bar z)}.   \tag{TA.21}
\]
One may alternatively use uniqueness of the normalized odd
theta period laws; the derivative at zero remains one.
Taking Taylor coefficients in (TA.21) is legitimate for
this specific continuous conjugation map. Conjugation
interchanges $w$ and $1-w$, so it preserves the two
nonvanishing clauses of the distinguished half-point.

Every embedding of $\mathbb C$ restricts on the quadratic
field $K$ either to the identity or to conjugation. In the
second case compose it with complex conjugation, apply the
fixed-$K$ result, and use (TA.21). TA1 identifies the new
conjugate lattice as $\bar\Lambda'$, since its two invariants
are the complex conjugates of those of $\Lambda'$. This
proves (TA.20) for every embedding of $\mathbb C$.

Finally let the original $\rho:M\hookrightarrow\mathbb C$
have only the literal coefficient assumptions in
\texttt{use1}. TA10 makes all even-theta coefficients
algebraic, and PA14 makes both invariants algebraic.
Extend the restriction of $\rho$ to
$M\cap\overline{\mathbb Q}$ to an automorphism of
$\mathbb C$, by the countable-field argument of TA5.
It agrees on both invariants and all coefficients to which
the statement applies. TA1 identifies the same
\texttt{lamConj}; therefore (TA.20) descends to the original
embedding. We have proved all three actual \texttt{use1}
clauses: algebraicity of every odd-theta rational translate,
algebraicity of every distinguished even theta coefficient,
and its prescribed conjugacy under the original $M,\rho$
and half-points. The quantifier order and coefficient-field
membership conditions are unchanged.

\paragraph{TA13: the updated analytic boundary.}
PA14, TA10 and PA15 now supply the primitive normalized
one-dimensional theta coefficients and all rational
translations; TA12 supplies the conjugacy conclusion of
\texttt{cmTheta\_taylorCoeff\_algebraic\_and\_conjugate}.
The generalized formal root argument already used in TL18
can therefore invoke these exact suppliers. PA4 also
supplies algebraic complex substitution, and PA10 supplies
the finite algebraic determinant correction $\kappa$.
This is a concrete reduction of the remaining vacuum
problem. It does not construct the actual
\texttt{StdOK}/four-copy vacuum identity, the normalized
indicator/covariance theorem, \texttt{Pins.iX}, or any of
the five later-stage existence assertions. Those remain
separately recorded targets. All proofs here are human
reconstructions with named classical inputs; no Lean
kernel replay or full self-contained RH/GRH proof is
claimed by this section.
